%%=============================================================================
%% Proof-of-concept implementatie
%%=============================================================================

\chapter{\IfLanguageName{dutch}{Proof-of-concept implementatie}{Proof-of-concept implementation}}%
\label{ch:poc}

In dit hoofdstuk wordt de proof-of-concept (PoC) besproken: de migratie van \texttt{AB\_APP\_GBP\_MAT\_BAF} naar PL/I,
toegepast volgens \emph{Component~2} (vertaalrichtlijnen) en \emph{Component~3} (best practices) van het
framework. Het hoofdstuk beschrijft eerst de gekozen modulaire opbouw, daarna de implementatiekeuzes per
submodule, en sluit af met een aantal afwijkingen ten opzichte van het CA~Gen-origineel die in de
validatiefase (Hoofdstuk~\ref{ch:validatie}) opnieuw aan bod komen.

\section{Architectuur van de PoC}
\label{sec:poc-architectuur}

De PL/I-implementatie volgt de bestaande PLAPO-conventies (zie \autoref{subsec:conventies}) en is opgesplitst
in 7 modules:
\begin{enumerate}
    \item \textbf{Voorloper} selecteert de populatie planbare materiaaleenheden volgens parameters.
    \item \textbf{Specifieke BAF MAIN-module} domeinspecifieke routering en parametrisatie; 
    orkestreert de aanroep van de materiaal- en bestellingmodules.
    \item \textbf{Generieke Materiaalmodule} produceert generieke materiaal-segmenten.
    \item \textbf{Specifieke Materiaalmodule} produceert specifieke materiaal-segmenten.
    \item \textbf{Generieke Bestellingmodule} produceert generieke bestellings-segmenten.
    \item \textbf{Specifieke Bestellingmodule} produceert specifieke bestellings-segmenten.
    \item \textbf{Plapo-manager} splitst verwerking op in logische eenheden en schrijft segmenten weg.
\end{enumerate}
